\section{Deployment Contract: Reliability, Loop, Observabilidad y Human Override}

La última slide didáctica reúne todo lo anterior en cuatro tipos de problemas de despliegue: confiabilidad, ciclos y deriva del plan, pruebas y benchmark, y despliegue real y observabilidad. No es un apéndice, sino el conjunto de condiciones de aceptación de la visión «Everyday AGI». Cuanto más se acerca el sistema al correo electrónico, al calendario, a los pagos y a las cuentas públicas del usuario, menos puede sustituirse la garantía operativa por la tasa de éxito promedio de una demo.

\lecturefigure{slide-58-key-issues-autonomous-agents.jpg}{Problemas clave de los Agent autónomos: confiabilidad, ciclos, pruebas y observabilidad}{00:46:09}

Esta página de síntesis se corresponde con la Everyday AGI de la apertura: aquella pregunta «¿qué pasaría si la inteligencia entrara en la vida cotidiana?», y esta responde «¿qué debe cumplirse antes de que entre?». La confiabilidad determina si se puede otorgar autorización; el loop control determina el costo y el límite de la pérdida de control; el testing determina si un cambio puede verificarse; la observabilidad determina si un incidente puede explicarse. Las cuatro no son parches de operación posteriores al lanzamiento, sino que deben diseñarse desde las etapas de agent architecture, recolección de datos y entorno de entrenamiento.

\subsection{Reliability: incluso con 99.9\% hay que preguntar por la frecuencia de la tarea y el costo del incidente}

Si la probabilidad de fallo de una sola acción de alto riesgo es $q$ y se ejecuta $N$ veces al año, la probabilidad de al menos un fallo es
\[
P(\ge 1\text{ failure})=1-(1-q)^N.
\]
Con $q=0.001$ y $N=1000$, esa probabilidad es de aproximadamente $63.2\%$. Por lo tanto, «99.9\% reliable» no es seguro por sí mismo; además hay que reducir las acciones automáticas de alto riesgo, añadir confirmaciones, limitar los permisos y garantizar que todo sea recuperable.

\teachervoice{00:46:09--00:50:00, el ponente dice que los pagos, la banca, el correo electrónico y las cuentas de redes sociales exigen una confiabilidad cercana al nivel de producción; el sistema no puede enviar contenido equivocado, realizar transacciones erróneas ni caer en ciclos. El docente enfatiza el audit trail, el monitoreo en línea, el human fallback y la toma de control por parte del usuario.}

\begin{warningbox}{La tasa de éxito promedio oculta la cola catastrófica}
Un informe de confiabilidad debe presentar a la vez la tasa de errores graves, la tasa de errores irreversibles, el uso de permisos sin confirmación, el ciclo más largo, el costo máximo por tarea y la tasa de toma de control humana. Reportar solo el éxito promedio de las tareas no demuestra que el sistema no provocará incidentes de baja frecuencia y alto daño.
\end{warningbox}

\subsection{Loop y plan divergence: dar al sistema un presupuesto y un estado detectable}

Un ciclo puede consistir en repetir la misma acción, o en un plan que se expande sin cesar sin avanzar. La detección puede combinar el hash del estado, los n-gram de acciones, la distancia al objetivo y el presupuesto:
\[
\text{loopRisk}_t=
\alpha\,\mathbf{1}[h(s_t)=h(s_{t-k})]
+\beta\,\text{actionRepeat}_t
+\gamma\,\text{noProgress}_t.
\]
Cuando el riesgo supera el umbral, el sistema debe detenerse, degradarse o solicitar intervención humana, en lugar de aumentar automáticamente el token budget.

\begin{lstlisting}[caption={Ciclo de ejecución con presupuesto, detección de repeticiones y toma de control humana}]
while not done:
    if spent_cost > cost_limit or elapsed > time_limit:
        return escalate("budget exceeded")
    if state_hash in recent_states:
        return escalate("loop detected")
    action = agent.next_action()
    if action.risk >= confirmation_threshold:
        action = wait_for_user_confirmation(action)
    execute(action)
    append_audit_log(action)
\end{lstlisting}

\subsection{Q\&A: de zero-shot a task-specific training}

Ante la pregunta «¿cómo pasar del 40\% a casi el 99\%?», el ponente señala que muchos frontier models operan en zero-shot frente a las agent interfaces: la distribución de preentrenamiento no incluye las páginas web ni los flujos de operación específicos. Mediante RL específico de la tarea, autocorrección y recolección repetida de trayectorias, un dominio estrecho puede elevarse; pero eso no equivale a que todas las tareas del mundo abierto se saturen al mismo tiempo.

\teachervoice{00:50:00--00:54:20, el ponente usa AgentQ para mostrar que el task-specific RL puede aumentar notablemente la exactitud en tareas de reservación, y a la vez reconoce que los sitios nuevos, las tareas nuevas y las restricciones nuevas siguen produciendo desplazamiento de distribución. La nota advierte: la saturación local y la confiabilidad general deben reportarse por separado.}

El ciclo virtuoso del despliegue puede escribirse como
\[
\begin{aligned}
\text{production traces}&\rightarrow\text{failure taxonomy}
\rightarrow\text{offline simulation}\\
&\rightarrow\text{training}\rightarrow\text{shadow eval}
\rightarrow\text{controlled rollout}.
\end{aligned}
\]
En cada paso deben conservarse las versiones y la provenance de los datos; solo así puede determinarse si una mejora proviene del modelo, del prompt, de las herramientas o de un cambio en el entorno.

\subsection{Hallucination y domain-specific regression suites}

La respuesta de ingeniería del ponente tiene dos niveles: un modelo base más potente reduce algunos errores; cada producto concreto sigue necesitando sus propias pruebas y evaluaciones. Para un agent vertical, debe mantenerse una regression suite formada por tareas reales, entradas extremas, límites de permisos, fallos de herramientas e incidentes históricos, y reproducirse a diario después de cualquier cambio en el modelo, el prompt, el schema de las herramientas o la policy.

\teachervoice{00:54:20--00:57:55, el ponente enfatiza construir alrededor de mil domain scenarios relevantes, revisar de forma continua los cambios de prompt y las regresiones, comparar distintos modelos y optimizar aún más con fine-tuning o reinforcement learning. Experiencia práctica: no tomar la actualización del modelo del proveedor como la aceptación propia.}

\begin{importantbox}{Un umbral de liberación entregable para un Agent}
\begin{enumerate}
  \item la strict offline suite no presenta regresiones graves;
  \item el shadow traffic no ejecuta efectos secundarios reales;
  \item se habilitan primero los usuarios canary y las herramientas de bajos permisos;
  \item las acciones de alto riesgo deben confirmarse y poder revertirse;
  \item el monitoreo cubre éxito, costo, ciclos, errores y toma de control humana;
  \item un incidente puede rastrearse hasta la versión del modelo, prompt, memory, tool o policy.
\end{enumerate}
\end{importantbox}

\subsection{Modelos pequeños, Manager--Worker y la prueba del mundo real}

La sesión de Q\&A discute smaller versus larger models. El ponente considera que un modelo pequeño entrenado con reasoning traces y RL puede ser más eficaz en tareas estrechas, y propone una posible arquitectura con un modelo grande como manager y modelos más pequeños como workers. Aquí debe respetarse el límite de la evidencia: el tamaño del modelo es solo una de las variables de enrutamiento; también cuentan la calidad, la latencia, el costo, la privacidad, el uso de herramientas y las necesidades de contexto.

\teachervoice{00:57:55--00:59:47, el ponente considera que se necesitan modelos «más inteligentes, no solo más grandes», y propone la posible combinación large manager + small workers; al mismo tiempo, toma el real-world task success como litmus test, en lugar de decidir según el número de parámetros.}

Si una tarea $x$ puede enrutarse a un modelo $m$, un objetivo sencillo es
\[
m^*(x)=\arg\max_m\left[
Q(m,x)-\lambda_L L(m,x)-\lambda_C C(m,x)-\lambda_R R(m,x)
\right],
\]
donde $Q$ es la calidad, $L$ la latencia, $C$ el costo y $R$ el riesgo. El enrutador debe elegir dinámicamente según la tarea, en lugar de fijar de manera permanente que «el manager debe ser el más grande y el worker el más pequeño».

\subsection{El problema de la memory no tiene una única respuesta de hardware}

La última pregunta compara la agent memory con la RAM, la ROM y el hard drive. El ponente no ofrece una arquitectura unificada, sino que enfatiza que cada aplicación es distinta y combina componentes distintos. Esta reserva es muy importante: algunas tareas solo necesitan un scratchpad de corto plazo, otras requieren un perfil de usuario editable y otras necesitan registros de eventos y una base de conocimiento; llamar «memory» a todo esto conduce a compartir información indebidamente y a una persistencia excesiva.

\teachervoice{00:59:47--01:00:57, ante la analogía RAM/ROM/disk, el ponente afirma explícitamente que no hay una straight answer; hay que buscar los modelos y componentes adecuados según la aplicación, como coding, chat o actions. Nota de clase: la memory architecture es una decisión de producto y de riesgo; más no siempre es mejor.}

\subsection{Resumen de la sección}

El contrato de despliegue convierte los indicadores de investigación del Agent en responsabilidades reales: la confiabilidad debe considerar la frecuencia de las tareas, los ciclos requieren detección y presupuesto, las pruebas deben cubrir la distribución del dominio, y en producción debe haber trace, monitoreo, auditoría y toma de control humana. El RL local puede mejorar tareas de dominio estrecho, pero no elimina el desplazamiento de distribución; tanto el tamaño del modelo como la estructura de la memory deben enrutarse según la tarea.
